% English translation of the verified Chinese baseline. Structural identifiers must remain synchronized.
\section{Conclusions}\label{sec:conclusion}

A unified framework has been developed for the multiphysics large-deformation dynamics of magnetically actuated flexible catheters, spanning a geometrically exact continuous model, joint Galerkin reduction, canonical Hamiltonian formulation, structured CHNN, and KAM stability analysis. The framework retains cross-sectional orientation, independent torsion, and rotary inertia and separates conservative action, passive dissipation, background-flow input, and control through mutually exclusive energy accounting. It thereby provides consistent state and power conventions for the continuous model, COMSOL data, and the finite-dimensional learning system.

The complex-representation-enhanced CHNN uses a real-valued physical Hamiltonian readout and independent branches for conservative correction, dissipation, and external ports. Statistical learning error is only one component of the KAM interface; complete certification additionally requires high-order or analytic error certificates, pre-existing physical perturbations, coordinate and truncation errors, symplectic-discretization error, nondegeneracy, and Diophantine conditions. The resulting analysis assigns distinct roles to continuous Hamiltonian KAM theory, KAM theory for symplectic maps, and a posteriori invariant-torus validation.

P0, P1, P1.5, and P1.6 provide low-dimensional or three-trajectory small-sample results. P1 shows a conditional CHNN advantage in the derivative-only comparison. P1.5 shows substantial residual-over-physics-only gains but no stable one-step-field advantage. P1.6 shows that all \(30/30\) constrained models satisfy the equilibrium conditions exactly while predictive errors deteriorate markedly. These results cannot be extrapolated into general CHNN superiority on the full P2 task, and different random seeds are not independent physical samples.

The \(0.5\%\) COMSOL--Galerkin Bridge perturbation-energy criterion is not met by all three trajectories, whereas the independent \(2\%\) force-v2 joint-L2 criterion is met by all three; the two findings have distinct interpretations. After same-trajectory comparison, error-source isolation, identification of the cotangent projection, comparison of solver-output semantics, and correction of fixed-basis initialization, the three-level study yields observed convergence orders of \(1.64186\)--\(2.00088\). In the \(1\,\mathrm{s}\) case, G2, pointwise G3, impulse G3, and peak-to-peak perturbation-energy drift are \(0.00146226\), \(0.000221162\), \(0.0000563606\), and \(0.00307057<0.005\), respectively. These results show that the current canonical data contract meets the data-generation requirements for subsequent multi-amplitude experiments.

Multi-amplitude statistics, the full 96-trajectory data set, the P2 model comparison, open-port analysis, high-fidelity discrepancy, and numerical KAM validation remain subjects of subsequent numerical work. This study has not shown BDF to be symplectic, has not obtained numerical KAM certification, an invariant torus, or an engineering safe domain, and has not established general CHNN superiority over Real HNN on the full P2 task. The theoretical guarantees are limited to the conservative backbone under the corresponding structural, regularity, nonresonance, and small-perturbation conditions. Open ports, contact switching, model discrepancy, and engineering safety must be evaluated jointly using the energy balance, port quantities, high-fidelity validation, and experimental data.
